\subsection{FREE FAMILIES. BASES}

Let A be a ring, T a set and consider the A-module \(\mathbf{A}_{s}^{\mathrm{(T)}}\) . By definition, it is the external direct sum of a family \((\mathbf{M}_{t})_{t\in\mathbb{T}}\) of A-modules all equal to \(A_{s}\) and for all \(t\in T\) there is a canonical injection \(j_{t}:A_{s}\to A_{s}^{\mathrm{(T)}}\) (no. 6). We write \(j_{t}(1) = e_{t}\) so that \(e_{t} = (\delta_{tt'})_{t' \in \mathbb{T}}\) , where \(\delta_{tt'}\) is equal to 0 if \(t' \neq t\) , to 1 if \(t' = t\) (``Kronecker symbol''; \((t, t') \mapsto \delta_{tt'}\) is just the characteristic function of the diagonal of \(T \times T\) ); every \(x = (\xi_{t})_{t \in T} \in A_{s}^{(T)}\) may then be written uniquely: \(x = \sum_{t \in T} \xi_{t} e_{t}\) . The mapping \(\phi: t \mapsto e_{t}\) of \(T\) into \(A_{s}^{(T)}\) is called canonical; it is injective if \(A\) is non-zero. We shall see that the ordered pair \((A_{s}^{(T)}, \phi)\) is a solution of a universal mapping problem (Set Theory, IV, § 3, no. 1).

PROPOSITION 17. For every A-module E and every mapping \(f: \mathbf{T} \to \mathbf{E}\) , there exists one and only one A-linear mapping \(g: \mathbf{A}_{s}^{(\mathrm{T})} \to \mathbf{E}\) such that \(f = g \circ \phi\) .

The condition \(f = g \circ \phi\) means that \(g(e_t) = f(t)\) for all \(t \in \mathbf{T}\) , which is equivalent to \(g(\xi e_t) = \xi f(t)\) for all \(\xi \in \mathbf{A}\) and all \(t \in \mathbf{T}\) and also means that \(g \circ j_t\) is the linear mapping \(\xi \mapsto \xi f(t)\) of \(\mathbf{A}_s\) into \(\mathbf{E}\) for all \(t \in \mathbf{T}\) ; the proposition is therefore a special case of no. 6, Proposition 6.

The linear mapping g is said to be determined by the family \((f(t))_{t\in\mathbb{T}}\) of elements of E; by definition

\[
g \left(\sum_ {t \in \mathrm{T}} \xi_ {t} e _ {t}\right) = \sum_ {t \in \mathrm{T}} \xi_ {t} f (t).\tag{43}
\]

The kernel \(\mathbf{R}\) of \(g\) is the set of \((\xi_t) \in \mathbf{A}_s^{(\mathrm{T})}\) such that \(\sum_{t} \xi_t f(t) = 0\) ; it is sometimes said that the module \(\mathbf{R}\) is the module of linear relations between the elements of the family \((f(t))_{t \in \mathbb{T}}\) . The exact sequence

\[
0 \longrightarrow \mathrm{R} \longrightarrow \mathrm{A} _ {s} ^ {(\mathrm{T})} \stackrel {\mathbf {g}} {\longrightarrow} \mathrm{E}\tag{44}
\]

is said to be determined by the family \((f(t))_{t\in\mathbb{T}}\) .

COROLLARY 1. Let T, T' be two sets, g: T → T' a mapping. Then there exists one and only one A-linear mapping f: A \(^{(T)}\) → A \(^{(T')}\) which renders commutative the diagram

\[
\begin{array}{c c c} \mathrm{T} & \stackrel {{g}} {{\longrightarrow}} & \mathrm{T} ^ {\prime} \\ \phi \Big \downarrow & & \Big \downarrow \phi^ {\prime} \\ \mathrm{A} ^ {(\mathrm{T})} & \stackrel {{f}} {{\longrightarrow}} & \mathrm{A} ^ {(\mathrm{T} ^ {\prime})} \end{array}
\]

where \(\phi\) and \(\phi'\) are the canonical mappings.

It suffices to apply Proposition 17 to the composite mapping \(\mathbf{T} \xrightarrow{\pmb{g}} \mathbf{T}' \xrightarrow{\phi'} \mathbf{A}^{(\mathbf{T}')}\) .

COROLLARY 2. For a family \((a_{t})_{t\in \mathbb{T}}\) of elements of an A-module \(\mathbf{E}\) to be a generating system of \(\mathbf{E}\) , it is necessary and sufficient that the linear mapping \(\mathbf{A}_{s}^{(T)}\to \mathbf{E}\) determined by this family be surjective.

This is just another way of expressing Proposition 9 of no. 7.

DEFINITION 10. A family \((a_{t})_{t\in \mathbf{T}}\) of elements of an A-module \(\mathbf{E}\) is called a free family (resp. a basis of \(\mathbf{E}\) ) if the linear mapping \(\mathbf{A}_{s}^{(\mathbf{T})}\to \mathbf{E}\) determined by this family is injective (resp. bijective). A module is called free if it has a basis.

In particular, a commutative group G is called free if G (written additively) is a free Z-module (cf.~I, § 7, no. 7).

Definition 10, together with Corollary 2 to Proposition 17, show that a basis of an A-module E is a free generating family of E. Every free family of elements of E is thus a basis of the submodule it generates.

By definition, the A-module \(\mathbf{A}_s^{(\mathrm{T})}\) is free and the family \((e_t)_{t\in \mathbb{T}}\) is a basis (called canonical) of this A-module. When \(\mathbf{A}\neq \{0\}\) , \(\mathbf{T}\) is often identified with the set of \(e_t\) by the canonical bijection \(t\mapsto e_t\) ; this amounts to writing \(\sum_{t\in \mathbb{T}}\xi_t.t\)

instead of \(\sum_{t\in T}\xi_{t}a_{t}\) for the elements of \(A_{s}^{(T)}\) . When this convention is adopted, the elements of \(A_{s}^{(T)}\) are called formal linear combinations (with coefficients in A) of the elements of T.

Definition 10 and Proposition 17 give immediately the following result:

COROLLARY 3. Let E be a free A-module, \((a_{t})_{t\in \mathbb{T}}\) a basis of E, F an A-module and \((b_{t})_{t\in \mathbb{T}}\) a family of elements of F. There exists one and only one linear mapping \(f:\mathbf{E}\to \mathbf{F}\) such that

\[
f (a _ {t}) = b _ {t} \text {   for   all   } t \in \mathrm{T}.\tag{45}
\]

For \(f\) to be injective (resp. surjective), it is necessary and sufficient that \((b_t)\) be a free family in \(\mathbf{F}\) (resp. a generating system of \(\mathbf{F}\) ).

When a family \((a_{t})_{t\in \mathbb{T}}\) is not free, it is called related. Definition 10 may also be expressed as follows: to say that the family \((a_{t})_{t\in \mathbb{T}}\) is free means that the relation \(\sum_{t\in \mathbb{T}}\lambda_t a_t = 0\) (where the family \((\lambda_t)\) is of finite support) implies \(\lambda_{t} = 0\) for all \(t\in \mathbb{T}\) ; to say that \((a_{t})_{t\in \mathbb{T}}\) is a basis of E means that every \(x\in \mathbb{E}\) can be written in one and only one way in the form \(x = \sum_{t\in \mathbb{T}}\xi_t a_t\) ; for all \(t\in \mathbb{T}\) , \(\xi_{t}\) is then called the component (or coordinate) of \(x\) of index \(t\) with respect to the basis \((a_{t})\) ; the mapping \(x\to \xi_t\) from E to \(\mathbf{A}_s\) is linear.

Suppose \(\mathbf{A} \neq \{0\}\) ; then, in an A-module \(\mathbf{E}\) , two elements of a free family \((a_t)_{t \in \mathbb{T}}\) whose indices are distinct are themselves distinct: for if \(a_{t'} = a_{t''}\) for \(t' \neq t''\) , then \(\sum_{t \in \mathbb{T}} \lambda_t a_t = 0\) with \(\lambda_{t'} = 1\) , \(\lambda_{t''} = -1\) and \(\lambda_t = 0\) for the elements of \(\mathbb{T}\) distinct from \(t'\) and \(t''\) . A subset \(\mathbb{S}\) of \(\mathbb{E}\) will be called a free subset (resp. a basis of \(\mathbb{E}\) ) if the family defined by the identity mapping of \(\mathbb{S}\) onto itself is free (resp. a basis of \(\mathbb{E}\) ); every family defined by a bijective mapping of an indexing set onto \(\mathbb{S}\) is then free (resp. a basis). The elements of a free subset of \(\mathbb{E}\) are also called linearly independent.

If a subset of E is not free, it is called related or a related system and its elements are said to be linearly dependent.

Every subset of a free subset is free; in particular, the empty subset is free and is a basis of the submodule \(\{0\}\) of E.

PROPOSITION 18. For a family \((a_{t})_{t\in \mathbb{T}}\) of a module \(\mathbf{E}\) to be free, it is necessary and sufficient that every finite subfamily of \((a_{t})_{t\in \mathbb{T}}\) be free.

This follows immediately from the definition.

Proposition 18 shows that the set of free subsets of E, ordered by inclusion, is inductive (Set Theory, III, § 2, no. 4); as it is non-empty (since \(\varnothing\) belongs to it), it has a maximal element \((a_{i})_{i\in I}\) by Zorn's Lemma. It follows (if \(\mathbf{A}\neq \{0\}\) ) that for all \(x\in \mathbf{E}\) there exist an element \(\mu \neq 0\) of A and a family \((\xi_{i})\) of element A such that \(\mu x = \sum_{i}\xi_{i}a_{i}\) (cf.~§ 7, no. 1).

PROPOSITION 19. Let E be an A-module, the direct sum of a family \((\mathbf{M}_{\lambda})_{\lambda \in L}\) of submodules. If, for each \(\lambda \in L\) , \(S_{\lambda}\) is a free subset (resp. generating set, basis) of \(M_{\lambda}\) , then \(S = \bigcup_{\lambda \in L} S_{\lambda}\) is a free subset (resp. generating set, basis) of E.

The proposition follows from the definitions and the relation \(\mathbf{A}_{s}^{(\mathbf{S})} = \bigoplus_{\lambda \in \mathbf{L}} \mathbf{A}_{s}^{(\mathbf{S}_{\lambda})}\) (associativity of direct sums, cf.~no. 6).

Remark. (1) By Definition 10, if \(\mathbf{A} \neq \{0\}\) and \((a_{t})_{t \in I}\) is a free family, no element \(a_{k}\) can be equal to a linear combination of the \(a_{t}\) of index \(t \neq k\) . But conversely, a family \((a_{t})\) satisfying this condition is not necessarily a free family. For example, let \(\mathbf{A}\) be an integral domain and \(a, b\) two distinct non-zero elements; in \(\mathbf{A}\) , considered as an A-module, \(a\) and \(b\) form a related system, since \((-b)a + ab = 0\) . But in general there does not exist an element \(x \in \mathbf{A}\) such that \(a = xb\) of \(b = xa\) (cf.~however § 7, no. 1, Remark).

An element x of a module E is called free if \(\{x\}\) is a free subset, that is if the relation \(\alpha x = 0\) implies \(\alpha = 0\) . Every element of a free subset is free and in particular 0 cannot belong to any free subset when \(A \neq \{0\}\) .

Remarks. (2) A free module can have elements \(\neq0\) which are not free: for example, the A-module \(A_{s}\) is free but the right divisors of zero in A are not free elements of \(A_{s}\) .

\begin{enumerate}
\def\labelenumi{(\arabic{enumi})}
\setcounter{enumi}{2}
\item
  In the additive group \(\mathbf{Z} / (n)\) ( \(n\) an integer \(\geqslant 2\) ) considered as a \(\mathbf{Z}\) -module, no element is free and a fortiori \(\mathbf{Z} / (n)\) is not a free module.
\item
  It can happen that every element \(\neq0\) of an A-module is free without the module being free. For example, the field Q of rational numbers is a Z-module with this property, for two elements \(\neq0\) of Q always form a related system and a basis of \(\mathbf{Q}\) could therefore only contain a single element \(a\) ; but the elements of \(\mathbf{Q}\) are not all of the form \(na\) with \(n \in \mathbf{Z}\) (cf.~VII, § 3).
\end{enumerate}

PROPOSITION 20. Every A-module E is isomorphic to a quotient module of a free A-module.

If T is a generating set of E, there exists a surjective linear mapping \(\mathbf{A}_{s}^{(\mathrm{T})}\to\mathbf{E}\) (Corollary 2 to Proposition 17), and if R is the kernel of this mapping, E is isomorphic to \(\mathbf{A}_{s}^{(\mathrm{T})}/\mathbf{R}\) .

In particular we may take \(\mathbf{T} = \mathbf{E}\) ; then there is a surjective linear mapping \(\mathbf{A}_s^{(\mathbf{E})} \to \mathbf{E}\) , called canonical.

In particular, to say that an A-module E is finitely generated (no. 7) means that it is isomorphic to a quotient of a free A-module with a finite basis or also that there exists an exact sequence of the form

\[
\mathrm{A} _ {s} ^ {n} \rightarrow \mathrm{E} \rightarrow 0 (n \text {   an   integer   } > 0).
\]

Note that if \(A \neq \{0\}\) every basis of a finitely generated free module E is necessarily finite, for if S is a finite generating system and B a basis of E, each element of S is a linear combination of a finite number of elements of B and if \(B'\) is the set of all the elements of B which figure thus in the expression for the elements of S, \(B'\) is finite and every \(x \in E\) is a linear combination of elements of \(B'\) , hence \(B' = B\) .

PROPOSITION 21. Every exact sequence of A-modules

\[
0 \longrightarrow \mathrm{G} \stackrel {g} {\longrightarrow} \mathrm{E} \stackrel {f} {\longrightarrow} \mathrm{F} \longrightarrow 0
\]

in which \(\mathbf{F}\) is a free A-module, splits (no. 9). To be precise, if \((b_{\lambda})_{\lambda \in \mathbf{L}}\) is a basis of \(\mathbf{F}\) and, for each \(\lambda \in \mathbf{L}\) , \(a_{\lambda}\) is an element of \(\mathbf{E}\) such that \(f(a_{\lambda}) = b_{\lambda}\) , the family \((a_{\lambda})_{\lambda \in \mathbf{L}}\) is free and generates a supplementary submodule to \(g(\mathbf{G})\) .

There exists one and only one linear mapping \(h: \mathbf{F} \to \mathbf{E}\) such that \(h(b_{\lambda}) = a_{\lambda}\) for all \(\lambda \in \mathbf{L}\) (Corollary 3 to Proposition 17). As \(h\) is a linear section associated with \(f\) , the proposition follows from I, § 4, no. 9, Proposition 15.

Remark. (5) Let \((a_{i})_{1\leqslant i\leqslant n}\) be a basis of an A-module E and let \((b_{i})_{1\leqslant i\leqslant n}\) be a family of elements of E given by the relations

\[
b _ {i} = \lambda_ {1 i} a _ {1} + \dots + \lambda_ {i i} a _ {i} (1 \leqslant i \leqslant n)\tag{46}
\]

where \(\lambda_{ii}\) is invertible in A; then \((b_i)_{1\leqslant i\leqslant n}\) is a basis of E. It suffices to argue by induction on \(n\) , the proposition being obvious for \(n = 1\) . If \(\mathbf{E}'\) is the submodule of E generated by the family \((a_i)_{1\leqslant i\leqslant n - 1}\) , it follows from the induction hypothesis that \((b_i)_{1\leqslant i\leqslant n - 1}\) is a basis of \(\mathbf{E}'\) ; on the other hand, it follows from (46) that if \(\mu b_n\in \mathbf{E}'\) with \(\mu \in \mathbf{A}\) , then also \(\mu \lambda_{nn}a_n\in \mathbf{E}'\) , whence \(\mu = 0\) since \(\lambda_{nn}\) is invertible. The family \((b_i)_{1\leqslant i\leqslant n}\) is thus free and, as

\[
a _ {n} = - \lambda_ {n n} ^ {- 1} \lambda_ {1 n} a _ {1} - \dots - \lambda_ {n n} ^ {- 1} \lambda_ {n - 1, n} a _ {n - 1} + \lambda_ {n n} ^ {- 1} b _ {n}
\]

it is seen that \((b_{i})_{1\leqslant i\leqslant n}\) is a generating system of E, which completes the proof. This result is easily generalized to a family \((a_{i})_{i\in I}\) whose indexing set I is well ordered.
% semantic-fragments: legacy-input-seam
\relax{}
